%% [langsci] alone: the \ea ... \z front-end of langsci-gb4e, with the dot
%% syntax NOT loaded.  Covers what the front-end has to get right on its
%% own -- the depth read off the nesting, the bracket judgment at every
%% level, \eal's head example with no text, and the deeper gloss stacks
%% the option brings with it.  langsci-mixed.tex covers the other half,
%% the two syntaxes in one document.
%%
%% Note for whoever edits this file: \ea in RUNNING TEXT opens an example.
%% Prose here must not mention it by name, which is why the sentences below
%% say "the front-end" where they would otherwise say the command.
\input{_preamble-langscionly}
\begin{document}
\ea\label{ls:one} LSONE the first example.
\z

\ea LSTWO the head of a nest.
  \ea LSSUBA the first sub-example.
  \ex[*]{LSSUBB judged through the bracket.}
    \ea LSROMANA the roman level.
    \z
  \ex LSSUBC back at the letters.
  \z
\ex LSTHREE a second top-level example, in the same block.
\z

\eal
  \ex LSLA under a head that has no text of its own.
  \ex LSLB the second letter.
\zl

\ea \gllll LSGOBJ zwei Zeilen\\
           LSGONE two lines\\
           LSGTWO deux lignes\\
           LSGTRI due righe\\
\glt LSTRANS the free translation.
\z

%% Justification.  This front-end deliberately does NOT set its examples
%% ragged right, where langsci-gb4e does; \ExRaggedRight asks for langsci's
%% behaviour.  Both need a line long enough to wrap before they can be told
%% apart, which is what the padding in these two is for -- they carry the
%% same text so that the only difference between them is the switch.
\ea LSJUST justified by default, and carrying enough text to wrap over
several lines, because one line reaches no margin at all and two might
reach it by accident: what tells the two settings apart is not whether
some line falls short of the margin, but whether the full lines all end
at the same place, which is what justification means and what ragged
setting gives up.
\z

{\ExRaggedRight
\ea LSRAG justified by default, and carrying enough text to wrap over
several lines, because one line reaches no margin at all and two might
reach it by accident: what tells the two settings apart is not whether
some line falls short of the margin, but whether the full lines all end
at the same place, which is what justification means and what ragged
setting gives up.
\z
}

LSREFS \ref{ls:one} and \Last.
\end{document}
